\documentclass[12pt,a4paper]{article}
\usepackage[a4paper,margin=18mm]{geometry}
\usepackage[T2A]{fontenc}
\usepackage[utf8]{inputenc}
\usepackage[russian]{babel}
\usepackage{amsmath,amssymb,mathtools}
\usepackage{array,booktabs,longtable}
\usepackage{xcolor}
\usepackage{hyperref}
\usepackage{tikz}

\hypersetup{
  colorlinks=true,
  linkcolor=blue!55!black,
  urlcolor=blue!55!black
}
\setlength{\parindent}{0pt}
\setlength{\parskip}{0.55em}
\setlength{\tabcolsep}{2pt}
\renewcommand{\arraystretch}{1.25}

\begin{document}

\begin{titlepage}
\thispagestyle{empty}
\begin{center}
\vspace*{0.24\textheight}
{\LARGE\bfseries Ревью домашней работы\par}
\vspace{2.2cm}
{\large Автор: Лёвин Артём Александрович\par}
\vspace{0.8cm}
{\large 8 сентября 2026 года\par}
\vfill
\begin{tikzpicture}
\draw[blue!45!black, line width=0.8pt]
  (0,0) .. controls (2.0,0.65) and (4.2,-0.65) .. (6.2,0)
        .. controls (8.0,0.55) and (10.0,-0.45) .. (12.0,0.05);
\end{tikzpicture}
\end{center}
\end{titlepage}

\newpage
\section*{Сводная таблица}
\small
\begin{longtable}{p{8mm}p{30mm}p{48mm}p{31mm}p{37mm}}
\toprule
\textbf{№} & \textbf{Тема} & \textbf{Суть ошибки} & \textbf{Ваш ответ} & \textbf{Правильный ответ} \\
\midrule
\endfirsthead
\toprule
\textbf{№} & \textbf{Тема} & \textbf{Суть ошибки} & \textbf{Ваш ответ} & \textbf{Правильный ответ} \\
\midrule
\endhead
\hyperref[task:6]{6} & Периодическая десятичная дробь & Не указана непериодическая часть десятичной записи, хотя это прямо требуется условием. & $0{,}21(6)$; период $6$ & $0{,}21(6)$; до периода: $21$; период $6$ \\
\addlinespace
\hyperref[task:7]{7} & Вынесение общего множителя; делимость & В цепочке равенств потерян множитель $41$; затем выполнено прямое умножение, которого условие просит избежать. & $82$ & $82$ \\
\addlinespace
\hyperref[task:10]{10} & Раскрытие скобок; приведение подобных; делимость & При перегруппировке слагаемое $+3b$ записано как $-3b$. Итог случайно получился верным, но промежуточное равенство неверно. & $3(a-b)$; кратно $3$ & $3(a-b)$; при целых $a,b$ делится на $3$ \\
\bottomrule
\end{longtable}
\normalsize

\newpage
\thispagestyle{empty}
\vspace*{\fill}
\begin{center}
{\LARGE\bfseries Разбор ошибок}
\end{center}
\vspace*{\fill}
\newpage

\phantomsection\label{task:6}
\section*{Задание 6}

\textbf{Текст задания.}
Для дроби
\[
\frac{13}{60}
\]
найдите десятичную запись, укажите непериодическую часть и период.

\textbf{Ваше решение.}
Выполнено деление и записано
\[
\frac{13}{60}=0{,}21(6).
\]
После этого указано: «Период $=6$».

\textbf{Ваш ответ.}
\[
0{,}21(6), \qquad \text{период }6.
\]
Непериодическая часть отдельно не указана.

\textcolor{red}{Ошибка:} не выполнено одно из обязательных требований условия --- не названа непериодическая часть периодической десятичной дроби.

\textbf{Где именно возникает ошибка.}
До записи
\[
\frac{13}{60}=0{,}21(6)
\]
ход решения верный. Также верно определён период $6$.
Неполнота возникает в самом конце: после определения периода нужно было ещё явно указать непериодическую часть.

\textbf{Почему это неверно.}
Запись
\[
0{,}21(6)=0{,}216666\ldots
\]
состоит из двух разных частей после запятой.
Цифры $2$ и $1$ появляются до начала бесконечного повторения; они образуют \emph{непериодическую часть} $21$.
Далее бесконечно повторяется одна цифра $6$; она образует период.
Поэтому одного сообщения «период $6$» недостаточно: условие требует назвать оба элемента структуры десятичной записи.

\textcolor{blue!60!black}{Ключевая идея:} в записи конечнопериодической дроби сначала выделяют цифры до первого повторяющегося блока, а затем сам минимальный повторяющийся блок.

\textbf{Исправление от последнего правильного шага.}
Из уже полученной правильной записи
\[
\frac{13}{60}=0{,}21(6)=0{,}216666\ldots
\]
сразу читаем:
\[
\text{непериодическая часть }=21,
\qquad
\text{период }=6.
\]

\textbf{Полное правильное решение.}
Разделим $13$ на $60$.
Поскольку $13<60$, целая часть равна нулю.
После добавления нуля получаем $130$: первая цифра после запятой равна $2$, так как
\[
130=2\cdot 60+10.
\]
Затем добавляем ноль к остатку $10$:
\[
100=1\cdot60+40,
\]
поэтому следующая цифра равна $1$.
Снова добавляем ноль:
\[
400=6\cdot60+40.
\]
Остаток снова равен $40$. Значит, дальше ситуация полностью повторяется: каждый следующий шаг снова даёт цифру $6$ и остаток $40$.
Следовательно,
\[
\frac{13}{60}=0{,}216666\ldots=0{,}21(6).
\]
До начала повторения стоят цифры $21$, поэтому непериодическая часть равна $21$. Повторяется цифра $6$, поэтому период равен $6$.

\textbf{Проверка.}
Можно проверить разложением дроби:
\[
\frac{13}{60}=\frac{12}{60}+\frac{1}{60}
=0{,}2+0{,}016666\ldots
=0{,}216666\ldots.
\]
Это совпадает с записью $0{,}21(6)$.

\textcolor{teal!60!black}{Итог:} десятичная запись --- $0{,}21(6)$; непериодическая часть --- $21$; период --- $6$.

\textbf{Задача на закрепление.}
Для дроби $\dfrac{17}{60}$ найдите десятичную запись, укажите непериодическую часть и период.

\newpage
\phantomsection\label{task:7}
\section*{Задание 7}

\textbf{Текст задания.}
Не выполняя прямого умножения больших чисел, докажите, что
\[
73\cdot41-39\cdot41
\]
делится на $17$. Найдите частное от деления полученного числа на $17$.

\textbf{Ваше решение.}
В решении записана цепочка, в которой встречается переход
\[
73\cdot41-39\cdot41=73-39=34,
\]
после чего найдено
\[
34:17=2.
\]
Далее записано корректное вынесение общего множителя
\[
41(73-39)=41\cdot34,
\]
но затем выполнено прямое умножение
\[
41\cdot34=1394,
\]
и деление $1394:17=82$.

\textbf{Ваш ответ.}
\[
82.
\]

\textcolor{red}{Ошибка:} в первом преобразовании потерян общий множитель $41$; кроме того, дальнейшее вычисление $41\cdot34=1394$ противоречит требованию решить задачу без прямого умножения больших чисел.

\textbf{Где именно возникает ошибка.}
Исходное выражение записано верно:
\[
73\cdot41-39\cdot41.
\]
Первый достоверно неверный переход:
\[
73\cdot41-39\cdot41=73-39.
\]
Такое равенство неверно, потому что множитель $41$ нельзя просто удалить из обоих слагаемых.
Позже в решении появляется правильная конструкция $41(73-39)$, поэтому эту часть рассуждения можно сохранить.

\textbf{Почему этот шаг неверен.}
Распределительный закон имеет вид
\[
ac-bc=c(a-b).
\]
Он позволяет \emph{вынести} общий множитель $c$ за скобки, но не позволяет его отбросить.
В данном случае
\[
73\cdot41-39\cdot41=41(73-39),
\]
а не $73-39$.
Если отбросить множитель $41$, значение выражения изменится в $41$ раз.

Кроме того, условие специально просит не выполнять прямое умножение больших чисел. После вынесения общего множителя это умножение вообще не нужно: достаточно заметить, что
\[
73-39=34=2\cdot17.
\]
Тогда множитель $17$ виден непосредственно.

\textcolor{blue!60!black}{Ключевая идея:} сначала вынести общий множитель $41$, затем представить разность $73-39$ как число, кратное $17$.

\textbf{Исправление от последнего правильного шага.}
Используем правильную часть Вашего решения:
\[
41(73-39).
\]
Вычисляем только небольшую разность:
\[
73-39=34=2\cdot17.
\]
Тогда
\[
41(73-39)=41\cdot34=41\cdot2\cdot17.
\]
Следовательно, исходное выражение делится на $17$, а частное равно
\[
41\cdot2=82.
\]

\textbf{Полное правильное решение.}
Вынесем общий множитель $41$:
\[
73\cdot41-39\cdot41=41(73-39).
\]
В скобках:
\[
73-39=34=2\cdot17.
\]
Поэтому
\[
41(73-39)=41\cdot2\cdot17=17\cdot(2\cdot41).
\]
Выражение представлено в виде $17$ умножить на целое число, значит, оно делится на $17$.
Частное равно
\[
2\cdot41=82.
\]
Прямое вычисление произведения $41\cdot34$ не потребовалось.

\textbf{Проверка.}
Проверим символически, не вычисляя большое произведение:
\[
17\cdot82=17\cdot(2\cdot41)=(2\cdot17)\cdot41=34\cdot41.
\]
С другой стороны,
\[
73\cdot41-39\cdot41=(73-39)\cdot41=34\cdot41.
\]
Обе записи совпадают, поэтому частное $82$ найдено верно.

\textcolor{teal!60!black}{Итог:} правильное преобразование --- $41(73-39)$; делимость следует из равенства $73-39=34=2\cdot17$; частное равно $82$.

\textbf{Задача на закрепление.}
Не выполняя прямого умножения больших чисел, докажите, что $68\cdot37-34\cdot37$ делится на $17$. Найдите частное от деления полученного числа на $17$.

\newpage
\phantomsection\label{task:10}
\section*{Задание 10}

\textbf{Текст задания.}
Упростите выражение
\[
0{,}6(15a-10b)-3(2a-b).
\]
Затем докажите, что при целых $a$ и $b$ значение выражения делится на $3$.

\textbf{Ваше решение.}
Сначала скобки раскрыты верно:
\[
0{,}6(15a-10b)-3(2a-b)=9a-6b-6a+3b.
\]
Далее записано
\[
9a-6a-6b-3b,
\]
после чего получено
\[
3a-3b=3(a-b)
\]
и сделан вывод, что выражение кратно $3$.

\textbf{Ваш ответ.}
\[
3(a-b), \qquad \text{выражение кратно }3.
\]

\textcolor{red}{Ошибка:} при перегруппировке слагаемых знак у $3b$ изменён с плюса на минус.

\textbf{Где именно возникает ошибка.}
Последний полностью правильный шаг:
\[
9a-6b-6a+3b.
\]
Первый неправильный шаг:
\[
9a-6b-6a+3b=9a-6a-6b-3b.
\]
При перестановке слагаемых $+3b$ должно остаться $+3b$.

\textbf{Почему этот шаг неверен.}
Перестановка слагаемых не меняет их знаков. Выражение
\[
9a-6b-6a+3b
\]
можно переписать как
\[
9a-6a-6b+3b,
\]
но не как $9a-6a-6b-3b$.

Если бы промежуточная строка
\[
9a-6a-6b-3b
\]
была верной, то приведение подобных дало бы
\[
3a-9b,
\]
а не $3a-3b$.
Следовательно, верный конечный результат получился только потому, что на следующем шаге фактически снова был использован исходный знак $+3b$. Такая компенсация не делает неверное промежуточное равенство допустимым.

\textcolor{blue!60!black}{Ключевая идея:} при раскрытии скобок и перегруппировке нужно переносить каждое слагаемое вместе с его знаком; затем отдельно объединять коэффициенты при $a$ и при $b$.

\textbf{Исправление от последнего правильного шага.}
Начинаем с последней правильной строки:
\[
9a-6b-6a+3b.
\]
Группируем слагаемые без изменения знаков:
\[
(9a-6a)+(-6b+3b).
\]
Тогда
\[
3a-3b=3(a-b).
\]

\textbf{Полное правильное решение.}
Раскроем первые скобки:
\[
0{,}6(15a-10b)=0{,}6\cdot15a-0{,}6\cdot10b=9a-6b.
\]
Раскроем вторые скобки с учётом минуса перед множителем $3$:
\[
-3(2a-b)=-6a+3b.
\]
Поэтому
\[
0{,}6(15a-10b)-3(2a-b)=9a-6b-6a+3b.
\]
Приводим подобные слагаемые:
\[
9a-6a=3a,
\qquad
-6b+3b=-3b.
\]
Получаем
\[
3a-3b=3(a-b).
\]
Теперь докажем делимость. Если $a$ и $b$ --- целые числа, то их разность $a-b$ тоже является целым числом. Следовательно, выражение
\[
3(a-b)
\]
представляет собой произведение числа $3$ на целое число. Значит, значение исходного выражения делится на $3$ при любых целых $a$ и $b$.

\textbf{Проверка.}
Проверим алгебраическое преобразование обратным раскрытием скобок:
\[
3(a-b)=3a-3b.
\]
А из исходного выражения после правильного раскрытия скобок получаем
\[
9a-6b-6a+3b=(9a-6a)+(-6b+3b)=3a-3b.
\]
Результаты совпадают.
Для дополнительной числовой проверки возьмём, например, $a=2$, $b=-1$:
\[
0{,}6(15\cdot2-10\cdot(-1))-3(2\cdot2-(-1))=24-15=9,
\]
а
\[
3(a-b)=3(2-(-1))=9.
\]

\textcolor{teal!60!black}{Итог:} правильное упрощение --- $3(a-b)$; при целых $a,b$ число $a-b$ целое, поэтому $3(a-b)$ делится на $3$.

\textbf{Задача на закрепление.}
Упростите выражение $0{,}5(18x-12y)-3(2x-y)$. Затем докажите, что при целых $x$ и $y$ значение выражения делится на $3$.

\end{document}
